% ECONOMICS_PROBLEM: {"id": "D44-102", "title": "Revenue-optimal multi-item, multi-bidder auctions", "jel_code": "D44", "source_id": "OP-0171"}
% !TeX program = xelatex
\documentclass[11pt,a4paper]{article}
\usepackage[margin=23mm]{geometry}
\usepackage{fontspec}
\setmainfont{Latin Modern Roman}
\usepackage{amsmath,amssymb,mathtools}
\usepackage{xcolor,enumitem,booktabs,longtable,array,xurl,hyperref}
\definecolor{muted}{HTML}{53636C}
\hypersetup{colorlinks=true,urlcolor=blue,linkcolor=blue}
\setlength{\parindent}{0pt}
\setlength{\parskip}{5pt}
\newcommand{\R}{\mathbb R}
\newcommand{\E}{\mathbb E}
\newcommand{\N}{\mathbb N}
\newcommand{\ind}{\mathbf 1}
\newcommand{\Var}{\operatorname{Var}}
\newcommand{\Cov}{\operatorname{Cov}}
\newcommand{\supp}{\operatorname{supp}}
\DeclareMathOperator*{\argmin}{arg\,min}
\DeclareMathOperator*{\argmax}{arg\,max}
\newcommand{\entrymeta}[3]{{\sffamily\small\color{muted}\textbf{\texttt{#1}}\quad|\quad #2\par #3\par}}
\newcommand{\Problemref}[1]{\texttt{#1}}
\newcommand{\JELref}[2]{\texttt{#2} (JEL \texttt{#1})}
\begin{document}
\section*{Revenue-optimal multi-item, multi-bidder auctions}
\textbf{Problem ID:} \texttt{D44-102}\quad \textbf{JEL:} \texttt{D44}\par
% BEGIN ECONOMICS STATEMENT
\label{problem:OP-0171}
% GAME_THEORY_PROBLEM: {"local_id": "GT-041", "original_number": 41, "kind": "R", "entry_kind": "statement_only", "published": false, "review_required": true, "category": "game_theory"}
\label{game:GT-041}
{\sffamily\small\color{muted}
\textbf{GT-041} | Mechanism design and auctions | \textbf{R}: rigorous formulation of a published characterization agenda.
\par}
\subsection*{Definitions and mathematical statement}
Fix $n,m\geq2$. Bidder $i$ has additive values $v_i\in V_i=[0,1]^m$ drawn independently across bidders from distributions $F_i$ with continuous densities. Let $F=\bigotimes_i F_i$. A direct randomized mechanism has measurable expected allocations $x_{ij}(v)\in[0,1]$ and integrable expected payments $p_i(v)$, with $\sum_i x_{ij}(v)\leq1$ for every $j,v$. Write
\[
 U_i(v_i;z_i,v_{-i})=\sum_j v_{ij}x_{ij}(z_i,v_{-i})-p_i(z_i,v_{-i}).
\]
Dominant-strategy incentive compatibility in expected utility requires $U_i(v_i;v_i,v_{-i})\geq U_i(v_i;z_i,v_{-i})$ for every $i,v_i,z_i,v_{-i}$; typewise individual rationality requires $U_i(v_i;v_i,v_{-i})\geq0$ at each realized type profile, in expectation over the mechanism's internal randomization. This IR convention does not require nonnegative utility for every realized lottery outcome. Let $\mathcal M_D(F)$ contain these mechanisms with $p_i(v)\geq0$. Determine the value, optimizer structure, and verifiable necessary-and-sufficient optimality conditions for
\[
 R_D(F)=\sup_{(x,p)\in\mathcal M_D(F)}\int \sum_i p_i(v)\,dF(v).
\]
For every input $F$ and $\varepsilon>0$, a constructive answer should specify a feasible mechanism of revenue at least $R_D(F)-\varepsilon$ and a matching upper bound. An efficiency claim additionally requires a finite representation of $F$ and an explicit accuracy-dependent running-time bound.
\subsection*{Short English statement}
Characterize the best truthful revenue for heterogeneous items and several privately informed additive bidders.
\subsection*{Variants and scope boundaries}
The Bayesian version replaces the dominant-strategy inequalities by their expectations over $v_{-i}\sim F_{-i}$ and typewise IR by interim IR. It defines a separate value $R_B(F)$; the two feasible sets must not be conflated. Universal truthfulness, which requires every realized deterministic mechanism to be truthful, is stronger than the expected-utility condition used here.
\subsection*{Source relationship and status}
The source explicitly describes the general characterization as unresolved and studies DSIC dual certificates. This entry supplies a bounded additive model and a precise optimization target; it is an agenda formulation, not a universal conjecture of efficient computability. The supremum convention does not presume attainment.
\subsection*{References}
\begingroup\footnotesize\raggedright
\textbf{[S1]} Yanchen Jiang, David C. Parkes, and Tonghan Wang (2026). \emph{Duality for Optimal Multi-Item, Multi-Bidder Auction Design: Revenue Certificates through Deep Learning}. arXiv:2606.10112v1. \textit{Verified locator:} abstract, first sentence and DSIC/discretization paragraphs. \url{https://arxiv.org/abs/2606.10112v1}
\par\endgroup

\subsection*{Common conventions: Source mathematical conventions}

\textbf{Finite games.} For a finite set $A$, $\Delta(A)$ is its probability
simplex. Players choose independent mixed strategies in Nash equilibrium;
correlation is allowed only when explicitly part of the solution concept.
In a game with payoff functions $u_i$ and strategy profile $x$, an additive
$\varepsilon$ Nash equilibrium satisfies
\[
 u_i(x)\geq u_i(x_i',x_{-i})-\varepsilon
 \quad\text{for every player $i$ and every admissible deviation $x_i'$}.
\]
For numerical approximation, payoffs are normalized as locally specified.
An $\varepsilon$ payoff residual is distinct from distance at most
$\varepsilon$ to the set of exact equilibria. Point convergence, convergence
of empirical play, and convergence in distance to an equilibrium set are
also distinct.

\textbf{Computational representations.} Unless stated otherwise, a complexity
question takes rational data with binary encoding; polynomial time is measured
in total bit length. A fixed-accuracy polynomial algorithm is not necessarily
polynomial in $1/\varepsilon$, and neither is necessarily polynomial in
$\log(1/\varepsilon)$. Succinct games and value, demand, or sampling oracles
must declare their interfaces. Exact real solutions require an explicit
representation; an unspecified list of arbitrary real numbers is not a
finite computational output. A claimed algorithmic barrier is meaningful
only for its specified model and reductions.

\textbf{Dynamic games.} Histories, information, observation, and allowable
randomization are part of the model. A stationary strategy depends only on
the current state; a positional strategy in a deterministic graph game
chooses one action at each controlled vertex. Nash equilibrium at an initial
state and equilibrium after every history are different requirements.
An expectation of a pathwise $\liminf$ payoff, a $\liminf$ of expected averages,
a limiting expected average, and a uniform finite-horizon equilibrium are
different criteria. None is silently substituted for another.

\textbf{Allocation and mechanisms.} A complete allocation assigns every item
exactly once unless otherwise stated. Zero-valued goods, negative values,
capacity constraints, feasibility, lotteries, and copies of goods affect
fairness definitions and are specified locally. Dominant-strategy incentive
compatibility, Bayesian incentive compatibility, universal truthfulness,
and truthfulness in expectation impose different inequalities. Whenever
an objective ratio has zero denominator, its convention is stated or those
instances are excluded.

\textbf{Infinite-dimensional models.} Expectations, integrals, stochastic
equations, and optimization objectives require their local regularity,
integrability, filtrations, and admissible domains. A backward--forward
mean-field system is a boundary-value problem; it is not an initial-value
semiflow without an additional construction. Long-time, large-population,
and vanishing-noise limits require an order of limits and a convergence
topology. If a minimum need not be attained, the objective is its infimum
or supremum and the approximation requirement is stated separately.
% END ECONOMICS STATEMENT
\end{document}
